\subsection{শিক্ষকের তথ্য থেকে নির্বাচন ও হিসাব}

নিচের সারণিতে MSc হলো \textbf{Master of Science} এবং MA হলো \textbf{Master of Arts} ডিগ্রির সংক্ষিপ্ত রূপ।

\begin{flushleft}\begin{minipage}{\linewidth}\raggedright

\small

\begin{tabular}{|l|l|l|l|l|}
\hline

\chSixHeader
TeacherID & Name & Subject & Qualification & Salary\\
\hline
T01 & Salma & ICT & MSc & 55000\\\hline
T02 & Rahat & Bangla & MA & 48000\\\hline
T03 & Farid & Physics & MSc & 62000\\\hline
T04 & Nabila & ICT & MSc & 58000\\\hline
T05 & Karim & English & MA & 45000\\
\hline

\end{tabular}

\captionof{table}{Teacher Table}

\end{minipage}\end{flushleft}


\textbf{অনুশীলন:}\begin{enumerate}
\item ICT বিষয়ের শিক্ষকদের তথ্য দেখাও।
\item যেসব শিক্ষকের Salary 50000-এর বেশি, তাঁদের তথ্য দেখাও।
\item Salary-এর সর্বোচ্চ মান বের করো।
\item শিক্ষকদের গড় Salary বের করো।
\end{enumerate}
\subsection{খেলোয়াড়ের তথ্য বিশ্লেষণ}

\begin{flushleft}\begin{minipage}{\linewidth}\raggedright

\small

\begin{tabular}{|l|l|l|l|l|l|}
\hline

\chSixHeader
PlayerID & Name & Team & Game & Age & Score\\
\hline
P01 & Sakib & Tigers & Cricket & 24 & 135\\\hline
P02 & Rafi & Eagles & Football & 22 & 18\\\hline
P03 & Tamim & Tigers & Cricket & 25 & 98\\\hline
P04 & Nabil & Warriors & Cricket & 21 & 152\\\hline
P05 & Rima & Queens & Badminton & 20 & 75\\
\hline

\end{tabular}

\captionof{table}{Player Table}

\end{minipage}\end{flushleft}


\textbf{অনুশীলন:}\begin{enumerate}
\item Cricket খেলোয়াড়দের তথ্য দেখাও।
\item যেসব খেলোয়াড়ের Score 100-এর বেশি, তাঁদের তথ্য দেখাও।
\item Score-এর সর্বোচ্চ মান বের করো।
\item Score-এর গড় বের করো।
\end{enumerate}
\subsection{শহর অনুযায়ী গ্রাহকের তথ্য বিশ্লেষণ}

\begin{flushleft}\begin{minipage}{\linewidth}\raggedright

\small

\begin{tabular}{|l|l|l|l|}
\hline

\chSixHeader
CustomerID & Name & Phone & City\\
\hline
C01 & Asif & 01710000001 & Dhaka\\\hline
C02 & Nila & 01710000002 & Rajshahi\\\hline
C03 & Anika & 01710000003 & Dhaka\\\hline
C04 & Karim & 01710000004 & Khulna\\\hline
C05 & Arif & 01710000005 & Dhaka\\\hline
C06 & Mitu & 01710000006 & Rajshahi\\
\hline

\end{tabular}

\captionof{table}{Customer Table}

\end{minipage}\end{flushleft}


\textbf{অনুশীলন:}\begin{enumerate}
\item Dhaka শহরের গ্রাহকদের তথ্য দেখাও।
\item মোট গ্রাহকসংখ্যা বের করো।
\item যেসব গ্রাহকের Name A দিয়ে শুরু, তাঁদের তথ্য দেখাও।
\item যেসব City-তে একাধিক গ্রাহক আছেন, সেগুলো দেখাও।
\end{enumerate}
